\section{Values, scales, labels and missing observations}
Defaults are \key{min=0} and \key{max=100}. Drawing ranges require min<max; include zero
when comparing magnitudes through bar lengths. Decimal/scientific literals and macros expanding to them are accepted;
expressions, thousands separators and units inside input are not.
Scientific notation is available as an output format.

Keep min, max and width identical within each comparison column.
Zero is not missing. Empty input or uppercase NA denotes missing data for a single
observation; missing stack segments are rejected. Overflow warns and clips geometry
while preserving the original label. Use \key{overflow=error} for strict checking.

\begin{demo}{Raw values, calculated percentages and missing values}
\gradbar[max=200]{50}

\gradbar[max=200,value format=percent]{50}

\gradbar[max=200,unit={\,ms}]{50}

\gradbar{0}

\gradbar[missing text={N/A}]{NA}
\end{demo}
\key{unit} appends text. \key{value format=percent} computes $100v/\mathrm{max}$,
ignores unit, and requires \key{min=0}. For values already expressed as percentages,
use \key{unit={\%}}. \key{text} overrides the complete label without changing geometry.
\key{precision} accepts integers from 0 to 6. \key{number format} is fixed, grouped
or scientific. Missing labels do not acquire a unit or percent suffix.

\clearpage

\subsection{Scientific input and finite precision}
Decimal and e/E scientific literals share graphics, CSV, selection, ranks,
statistics and range calculations. Zero is valid; nonzero magnitudes must lie
between $10^{-100}$ and $10^{100}$, including explicit ranges and displayed
statistics. Each decimal side accepts up to 101 digits; the exponent accepts
up to three digits. Internal arithmetic has about 16 significant digits, so
additional digits do not guarantee distinct values. Locale separators,
expressions, input units, NaN and infinity are errors. Display precision does
not alter calculations or ranks.
\begin{demo}{Scientific input preserves the raw value}
\gradbar[min=0,max=1e-2,number format=scientific,precision=2]{1.25E-3}
\end{demo}
